% -------------------------------------------------------------------------
% WBS ID: RES-ML-TRN
% Deliverable: Training, Validation and Experiment Methodology
% Status: In progress
% Evidence status: ML training deferred by current hydrodynamic scope
% Review status: Not reviewed
% Last updated: 2026-07-19
% Dependencies: RES-ML-DAT, RES-ML-ALG, RES-ML-OBJ
% Downstream interfaces: RES-ML-GEN, RES-ML-SPEC
% -------------------------------------------------------------------------

\section{RES-ML-TRN --- Training, Validation and Experiment Methodology}
\label{sec:res-ml-trn}

\subsection{Purpose and Scope}
\label{sec:res-ml-trn-purpose}

This Deliverable records future training, grouped validation and
model-selection rules only. There is no active ML training in the
current hydrodynamic Research phase. If ML is reactivated later, the
simulation run shall be the indivisible grouping unit.

\subsection{Research Questions}
\label{sec:res-ml-trn-research-questions}

% TODO: State the questions that must be answered before the Deliverable
% can be accepted.

\subsection{Terminology and Definitions}
\label{sec:res-ml-trn-definitions}

% TODO: Define the technical terminology, symbols, quantities and
% classifications used in this Deliverable.

\subsection{Literature Evidence}
\label{sec:res-ml-trn-literature-evidence}

% TODO: Record evidence from peer-reviewed papers, authoritative datasets,
% standards, technical reports and primary documentation.
%
% Do not create a detached literature-summary list. Explain how each source
% supports, contradicts or limits a project decision.

\subsection{Governing Relations and Quantitative Definitions}
\label{sec:res-ml-trn-governing-relations}

% TODO: Record relevant equations, dimensionless groups, empirical models,
% extraction rules or quantitative definitions.
%
% Use align environments for displayed equations.
% Every displayed equation must have a unique \label{eq:...}.
% Do not place punctuation after displayed equations.

\subsection{Machine-Learning Role}
\label{sec:res-ml-trn-machine-learning-role}

% TODO: Complete this RES-ML Domain-specific subsection.

\subsection{Non-ML Baseline}
\label{sec:res-ml-trn-non-ml-baseline}

% TODO: Complete this RES-ML Domain-specific subsection.

\subsection{Input and Output Representation}
\label{sec:res-ml-trn-input-output-representation}

% TODO: Complete this RES-ML Domain-specific subsection.

\subsection{Dataset Requirements}
\label{sec:res-ml-trn-dataset-requirements}

% TODO: Complete this RES-ML Domain-specific subsection.

\subsection{Model or Architecture Candidates}
\label{sec:res-ml-trn-model-architecture-candidates}

% TODO: Complete this RES-ML Domain-specific subsection.

\subsection{Training and Validation Method}
\label{sec:res-ml-trn-training-validation-method}

If ML is reactivated later, the previous secondary-event holdout
requirement shall be replaced with:
\begin{itemize}
  \item grouped splits by complete simulation;
  \item held-out geometry or intervention options;
  \item held-out condition ranges;
  \item structure-level holdout where sufficient structures exist;
  \item grouped or nested cross-validation;
  \item fixed seeds;
  \item non-ML baselines;
  \item independent final test cases.
\end{itemize}

Time steps or spatial samples from one simulation shall not be randomly divided
between future training and test partitions.

\subsection{Evaluation Metrics}
\label{sec:res-ml-trn-evaluation-metrics}

% TODO: Complete this RES-ML Domain-specific subsection.

\subsection{Generalisation and Uncertainty}
\label{sec:res-ml-trn-generalisation-uncertainty}

% TODO: Complete this RES-ML Domain-specific subsection.

\subsection{Simulation and Optimisation Integration}
\label{sec:res-ml-trn-simulation-optimisation-integration}

% TODO: Complete this RES-ML Domain-specific subsection.

\subsection{Candidate Approaches}
\label{sec:res-ml-trn-candidate-approaches}

% TODO: Define the credible candidate approaches considered.

\subsection{Critical Comparison}
\label{sec:res-ml-trn-critical-comparison}

% TODO: Compare candidates using physically and project-relevant criteria.
% Do not select an approach solely on popularity or implementation ease.

\subsection{Assumptions and Applicability}
\label{sec:res-ml-trn-assumptions}

% TODO: State assumptions and the conditions under which they are valid.

\subsection{Limitations and Failure Conditions}
\label{sec:res-ml-trn-limitations}

% TODO: State where the evidence, model or selected approach ceases to be
% reliable or applicable.

\subsection{Project-Specific Conclusions}
\label{sec:res-ml-trn-conclusions}

Future ML validation shall be grouped by engineering case and
geometry/intervention option, not by arbitrary sample rows. This
protects against leakage from repeated outputs of the same simulation
campaign.

\subsection{Selected Approach and Rationale}
\label{sec:res-ml-trn-selected-approach}

% TODO: Record the selected approach only after sufficient evidence exists.
% State the alternatives rejected and the reasons for rejection.

\subsection{Unresolved Decisions}
\label{sec:res-ml-trn-unresolved-decisions}

% TODO: Record unresolved questions, required evidence and decision owners.

\subsection{Requirements and Handoffs}
\label{sec:res-ml-trn-requirements}

% TODO: State requirements passed to other Research Domains, Software,
% verification activities or later execution work.

\subsection{Deliverable Completion Record}
\label{sec:res-ml-trn-completion-record}

% TODO: Record whether the Deliverable is:
%
% - Not started
% - In progress
% - Draft complete
% - Reviewed
% - Accepted
%
% Acceptance requires:
%
% - the research questions to be answered;
% - claims to be supported by appropriate evidence;
% - assumptions and limitations to be explicit;
% - the project-specific conclusion to be stated;
% - downstream requirements to be traceable;
% - unresolved decisions to be closed or formally deferred.
